在运动学中，角速度和角加速度描述转动怎样进行；在动力学中，还需要知道转动状态怎样受力而改变。
地球的自转、收拢双臂的滑冰者和高速旋转的陀螺，都涉及这一问题。相应的物理量是角动量，
改变角动量的是力矩。零外力矩保证角动量不变，却不保证任意刚体的角速度方向都不变。

\subsection{外积与转动方向}
内积把两个矢量联系成标量，适合描述投影、功和功率；外积又称叉积，适合描述有方向的转动。
对非零矢量 $\vect a,\vect b$，取它们同起点时的夹角 $0\le\theta\le\pi$，有
\begin{equation}
 |\vect a\times\vect b|=ab\sin\theta.
\end{equation}
其大小是以两矢量为邻边的平行四边形面积，方向垂直于该平面。右手四指沿小于或等于
$180^\circ$ 的转角从 $\vect a$ 弯向 $\vect b$，拇指所指就是外积方向。平行或反平行时
外积为零，不再给零矢量指定方向。图~\ref{ch09:fig:cross}中从向右的 $\vect a$ 转到右上的
$\vect b$ 为逆时针，故外积出纸面，以 $\odot$ 表示；$\otimes$ 表示入纸面。
\crossProductDiagram
外积满足反交换律和分配律，但一般不满足结合律：
\begin{align}
 \vect a\times\vect b&=-\vect b\times\vect a,\qquad \vect a\times\vect a=\vect0,\\
 (\vect a+\vect b)\times\vect c&=\vect a\times\vect c+\vect b\times\vect c,\\
 (\vect a\times\vect b)\times\vect c&=(\vect a\cdot\vect c)\vect b-(\vect b\cdot\vect c)\vect a,\\
 \vect a\times(\vect b\times\vect c)&=(\vect a\cdot\vect c)\vect b-(\vect a\cdot\vect b)\vect c.
\end{align}
在右手直角坐标系中，分量计算为
\begin{equation}
 \vect a\times\vect b=
 (a_yb_z-a_zb_y,\ a_zb_x-a_xb_z,\ a_xb_y-a_yb_x).
\end{equation}
例如纸面为 $xy$ 平面，向右、向上分别为正方向，$+z$ 就是出纸面。这个约定下
$(\vect r\times\vect p)_z=xp_y-yp_x$，正值对应逆时针趋势，负值对应顺时针趋势。
这是右手定则的代数形式，不能在一道题中途更换正方向。

位置、速度、动量属于极矢量，角速度、力矩和角动量属于轴矢量。以空间反演为例，
极矢量 $\vect r,\vect p$ 都变号，而 $(-\vect r)\times(-\vect p)=\vect r\times\vect p$ 不变。
因此外积的变换性质不同于普通位置矢量；例如轴矢量与极矢量的外积
$\vect v=\vect\omega\times\vect r$ 仍是极矢量。

\subsection{角动量、力矩与参考点}
先在一个惯性系内取固定点 $O$。质量恒定为 $m$ 的质点位置为 $\vect r=\overrightarrow{OP}$，
速度为 $\vect v$，动量为 $\vect p=m\vect v$。定义
\begin{equation}
 \vect L_O=\vect r\times\vect p,\qquad L_O=mrv\sin\varphi=p\,d,
\end{equation}
其中 $\varphi$ 是 $\vect r$ 与 $\vect v$ 的夹角，$d$ 是 $O$ 到动量作用线的垂距，
不是 $O$ 到质点的距离。图~\ref{ch09:fig:particle}把 $\vect r$ 平移到 $P$ 再量角，
两矢量的夹角为锐角；$p_y=0,y>0$，故 $L_z=-yp_x<0$，角动量入纸面。
\particleAngularMomentumDiagram
若质点沿过 $O$ 的直线径向运动，则 $d=0$、$L_O=0$；若以恒定速度沿不经过 $O$ 的直线运动，
$d$ 不变，角动量也不变。角动量必须附带参考点，同一动量对不同点的角动量一般不同。

由于 $O$ 固定，$\dd\vect r/\dd t=\vect v$。对定义求导，并利用牛顿第二定律：
\begin{align}
 \frac{\dd\vect L_O}{\dd t}
 &=\frac{\dd\vect r}{\dd t}\times m\vect v+
   \vect r\times\frac{\dd(m\vect v)}{\dd t}\\
 &=\vect v\times m\vect v+\vect r\times\vect F
 =\vect r\times\vect F.
\end{align}
第一项为零，因为两个矢量平行。定义对同一点的力矩 $\vect M_O=\vect r\times\vect F$，得
\begin{equation}
 \boxed{\frac{\dd\vect L_O}{\dd t}=\vect M_O}.\label{ch09:eq:angular-law}
\end{equation}
力矩大小是力乘以相应垂直力臂，平行于 $\vect r$ 的力分量不贡献力矩。合力非零，合力矩仍可为零。
角动量单位为 $\mathrm{kg\,m^2/s}$，力矩单位为 $\mathrm{N\,m}=\mathrm{kg\,m^2/s^2}$，
二者相差一个时间因子。积分形式为
\begin{equation}
 \vect L_O(t_2)-\vect L_O(t_1)=\int_{t_1}^{t_2}\vect M_O\,\dd t.
\end{equation}
右侧称为冲量矩。碰撞时外力瞬时很大不一定破坏近似守恒，但必须确认该外力的冲量矩可忽略，
或其作用线通过所取矩心；不能仅凭“时间短”就一概使用角动量守恒。

\subsection{质点系、质心分解及移动矩心}
对质量不变的质点系，设总质量为 $M$，总动量为 $\vect P$，质心位置、速度分别为
$\vect R_C,\vect V_C$。令 $\vect r_i=\vect R_C+\vect r_i'$、
$\vect v_i=\vect V_C+\vect v_i'$。质心定义给出
\begin{equation}
 \sum_i m_i\vect r_i'=\vect0,\qquad
 \sum_i m_i\vect v_i'=\frac{\dd}{\dd t}\sum_i m_i\vect r_i'=\vect0.
\end{equation}
逐项展开角动量求和，而不是把混合项略去：
\begin{align}
 \vect L_O&=\sum_i m_i(\vect R_C+\vect r_i')\times(\vect V_C+\vect v_i')\\
 &=M\vect R_C\times\vect V_C
 +\vect R_C\times\sum_i m_i\vect v_i'
 +\left(\sum_i m_i\vect r_i'\right)\times\vect V_C
 +\sum_i m_i\vect r_i'\times\vect v_i'\\
 &=\boxed{M\vect R_C\times\vect V_C+\vect L_C}.
\end{align}
第一项称轨道角动量，第二项是相对质心的内部角动量。以地球为例，对太阳取矩时，
公转与自转分别对应这两项。此分解不要求系统是刚体：粒子团、爆炸碎片和可形变物体同样适用。

设 $\vect F_{ij}$ 是 $j$ 对 $i$ 的内力，若 $\vect F_{ji}=-\vect F_{ij}$ 且内力沿两质点连线，
则一对内力的力矩和为
\begin{equation}
 \vect r_i\times\vect F_{ij}+\vect r_j\times\vect F_{ji}
 =(\vect r_i-\vect r_j)\times\vect F_{ij}=\vect0.
\end{equation}
因此总角动量的变化只由合外力矩决定。若其全部分量为零，则角动量矢量守恒；
若只有某一固定方向的力矩分量为零，则只能断定那个方向的角动量分量守恒。
轨道角动量的变化率是 $\vect R_C\times\vect F_{\rm ext}$，即把合外力作用于质心后
对 $O$ 的力矩；自转角动量的变化率是各外力对质心的合力矩。两项相加才是各外力对O的总力矩，
不能把总力矩全部算进轨道部分。

\textbf{换固定矩心。}若 $O,O'$ 都固定在同一惯性系中，$\vect a=\overrightarrow{OO'}$，则
\begin{align}
 \vect L_{O'}&=\vect L_O-\vect a\times\vect P,\\
 \vect M_{O'}&=\vect M_O-\vect a\times\vect F_{\rm ext}.
\end{align}
总动量为零时角动量与固定参考点选择无关；合外力为零时外力矩与固定参考点选择无关。
一般情形两者均随矩心改变。选支点是为消去支点力矩，并不意味着其他外力矩自动为零。

\textbf{矩心移动。}若 $O$ 以速度 $\vect v_O$ 移动，仍用惯性系中的绝对速度定义
$\vect L_O^{\rm abs}=\sum_i(\vect r_i-\vect R_O)\times m_i\vect v_i$，则
\begin{equation}
 \frac{\dd\vect L_O^{\rm abs}}{\dd t}
 =\sum_i(\vect v_i-\vect v_O)\times m_i\vect v_i+\vect M_O
 =\vect M_O-\vect v_O\times\vect P.
\end{equation}
所以即使矩心匀速运动，也不能不加条件地沿用固定点公式。
若改用相对速度定义 $\vect L_O^{\rm rel}$，并让参考轴只平移、不旋转，则
\begin{equation}
 \frac{\dd\vect L_O^{\rm rel}}{\dd t}
 =\vect M_O-\left[\sum_i m_i(\vect r_i-\vect R_O)\right]\times\vect a_O.
\end{equation}
两种公式对应不同的角动量定义，不应混用。特别取 $O=C$ 时，括号中的量为零。
在只随质心平移的坐标系里，也可把每个质点的惯性力写成 $-m_i\vect a_C$，其对质心的合力矩
\begin{equation}
 \sum_i\vect r_i'\times(-m_i\vect a_C)
 =-\left(\sum_i m_i\vect r_i'\right)\times\vect a_C=\vect0.
\end{equation}
这才是质心即使加速，仍有 $\dd\vect L_C/\dd t=\vect M_C^{\rm ext}$ 的原因。

\subsection{有心力、回归点与有效势}
力始终沿质点与固定力心的连线，称为有心力。它对力心的力矩为零，角动量守恒；
若 $\vect L\ne\vect0$，$\vect r,\vect v$ 始终垂直于固定的 $\vect L$，轨迹处于一个固定平面。
$L=0$ 对应径向运动，应单独处理。平面极坐标中取沿运动的方位角 $\varphi$，有
\begin{equation}
 \vect v=\dot r\,\uvec e_r+r\dot\varphi\,\uvec e_\varphi,
 \qquad L=mr^2\dot\varphi.
\end{equation}
窄扇形面积为 $\dd A=\tfrac12r^2\dd\varphi$，故面积速度
\begin{equation}
 \frac{\dd A}{\dd t}=\frac12r^2\dot\varphi=\frac{L}{2m}.
\end{equation}
这对任何固定中心力成立，不只适用于平方反比引力；面积矢量的直接推导见后面的面积速度例题。

质点到力心的距离达到极值的位置称回归点，此时 $\dot r=0$，速度只有横向分量。
若有两个回归点 $r_1,r_2$，便有 $r_1v_1=r_2v_2$；回归点不是合力为零的点。
图~\ref{ch09:fig:apsides}中两处速度都与日心连线垂直，近点半径小，速率反而大。
\apsidesDiagram
若中心力又是保守力，势能为 $U(r)$，把速度分解代入机械能：
\begin{align}
 E&=\frac12m\dot r^2+\frac12mr^2\dot\varphi^2+U(r)\\
 &=\frac12m\dot r^2+U_{\rm eff}(r),
 \qquad U_{\rm eff}(r)=U(r)+\frac{L^2}{2mr^2}.
\end{align}
有效势中的第二项不是新增的真实相互作用，而是角动量约束在径向方程中的体现。
可达范围必须有 $E\ge U_{\rm eff}$，回归点满足等号。圆轨道还要求
$U'_{\rm eff}(r_0)=0$；若 $U''_{\rm eff}(r_0)>0$，小径向扰动稳定，负值则不稳定，
二阶导数为零时不能仅凭二阶判据下结论。能量给速度大小，角动量把横向速度与半径联系起来，二者不能互相替代。

以 $M\gg m$ 的万有引力为例，$U=-GMm/r$。在回归点代入 $v=L/(mr)$，令 $u=1/r$：
\begin{equation}
 \frac{L^2}{2m}u^2-GMm\,u-E=0,
 \quad
 u_\pm=\frac{GMm^2}{L^2}
 \left(1\pm\sqrt{1+\frac{2EL^2}{G^2M^2m^3}}\right).
\end{equation}
这里已假定 $L\ne0$。允许能量须满足判别式非负：
$E\ge E_{\min}=-G^2M^2m^3/(2L^2)$。
当 $E_{\min}<E<0$ 时，两根均正，对应椭圆的近、远日点；$E=E_{\min}$ 时重根对应圆轨道。
$E=0$ 时一个根为零，即远端趋向无穷远；$E>0$ 时仅一个根为正，只存在一个有限距离的回归点。
图~\ref{ch09:fig:kepler}专画束缚椭圆情形，太阳位于焦点而非中心，两个半径都从太阳量起。
\keplerContextDiagram

\subsection{中心力轨道方程与开普勒定律}
进一步把时间消去，可以直接求轨道形状。设 $h=L/m=r^2\dot\varphi\ne0$，令 $u=1/r$，
撇号在本小节表示对 $\varphi$ 求导。由 $\dot\varphi=hu^2$，
\begin{equation}
 \dot r=-\frac{u'}{u^2}\dot\varphi=-hu',\qquad
 \ddot r=-h u''\dot\varphi=-h^2u^2u'',\qquad
 r\dot\varphi^2=h^2u^3.
\end{equation}
把它们代入径向牛顿方程 $m(\ddot r-r\dot\varphi^2)=F(r)$，得到比内方程（Binet equation）
\begin{equation}
 u''+u=-\frac{F(1/u)}{mh^2u^2}.
\end{equation}
引力 $F=-GMmu^2$ 使右侧成为常量 $GM/h^2$，故
\begin{equation}
 u=\frac{GM}{h^2}+A\cos\varphi+B\sin\varphi
 =\frac1p[1+e\cos(\varphi-\varphi_0)],
 \quad p=\frac{h^2}{GM},\quad e=p\sqrt{A^2+B^2}.
\end{equation}
其中 $p$ 称半通径，$e\ge0$ 是离心率，$\varphi_0$ 给出近点方位；于是
$r=p/[1+e\cos(\varphi-\varphi_0)]$。$0\le e<1$ 是椭圆，$e=1$ 是抛物线，
$e>1$ 是双曲线，圆是 $e=0$ 的特殊情形。各类均须限制 $r>0$ 的方位范围。
椭圆半长轴 $a$、半短轴 $b$、焦距 $c$ 满足
$c=ae=\sqrt{a^2-b^2}$，$p=a(1-e^2)=b^2/a$。

对椭圆，$r_{\rm p}=a(1-e)$、$r_{\rm a}=a(1+e)$。
把近点半径及 $L^2=GMm^2p$ 代入能量，逐步化简得
\begin{align}
 E&=\frac{GMm\,a(1-e^2)}{2a^2(1-e)^2}-\frac{GMm}{a(1-e)}\\
 &=\frac{GMm}{2a(1-e)}[(1+e)-2]
 =-\frac{GMm}{2a}.
\end{align}
一周扫过面积为 $\pi ab$，面积速度为 $h/2$，所以
\begin{equation}
 T=\frac{2\pi ab}{h},\qquad
 T^2=\frac{4\pi^2a^2b^2}{GMp}=\frac{4\pi^2a^3}{GM}.
\end{equation}
这分别给出轨道能量与开普勒第三定律。$a$ 是半长轴，不是近点距离；同一中心天体周围的
椭圆轨道满足 $T^2\propto a^3$，与面积速度恒定一起把几何、时间与动力学联系起来。

\subsection{转动惯量及两条移轴定理}
刚体各质元相对位置不变，绕固定轴转动时有共同角速度。对轴向角动量，有
\begin{equation}
 L_{\rm axis}=I_{\rm axis}\omega,\qquad
 I_{\rm axis}=\sum_i m_i r_{i\perp}^2=\int r_\perp^2\dd m.
\end{equation}
$r_\perp$ 是到轴的垂距。仅当该轴是相应的主惯量轴时，才可把整个角动量写成
$\vect L=I\vect\omega$；一般情况下还有垂直于该轴的角动量分量。
在三维中可写 $\vect L_C=\mathbf I_C\vect\omega$，$\mathbf I_C$ 称惯量张量，
其作用不一定只是乘一个数，也可能改变方向。后面用两质点刚体具体说明。

图~\ref{ch09:fig:inertia}列出八种均匀物体及指定转轴。薄圆筒的质量全在半径R处，故为 $MR^2$；
实心圆柱的一部分质量更靠近轴，故为 $MR^2/2$。形状、尺寸含义和轴位置必须同时核对。
\momentsContextDiagram
例如边长 $a,b$ 的矩形薄板，取过中心且垂直板面的轴。面密度为 $M/(ab)$，有
\begin{align}
 I&=\frac{M}{ab}\int_{-a/2}^{a/2}\int_{-b/2}^{b/2}(x^2+y^2)\,\dd y\,\dd x\\
 &=\frac{M}{ab}\left(b\frac{a^3}{12}+a\frac{b^3}{12}\right)
 =\frac{M(a^2+b^2)}{12}.
\end{align}
量纲为质量乘长度平方，$b\to0$ 时回到细杆过中心垂轴的结果。

设目标轴与过质心的轴平行，轴间距为d。质元到质心轴的横向矢量为 $\vect\rho_i$，
到目标轴的横向矢量为 $\vect\rho_i+\vect d$。展开平方得到平行轴定理：
\begin{align}
 I&=\sum_i m_i|\vect\rho_i+\vect d|^2\\
 &=I_C+2\vect d\cdot\sum_i m_i\vect\rho_i+Md^2
 =\boxed{I_C+Md^2}.
\end{align}
中间项由质心定义消失。这解释了同一细杆绕端点的惯量从 $ML^2/12$ 增加到 $ML^2/3$。
两轴必须平行；复合体各部分先移到同一轴再相加，挖孔则减去被移除部分关于同一轴的惯量。

若所有质元在同一薄平面内，x、y轴在板面内且互相垂直，z轴过同一点垂直板面，则
\begin{equation}
 I_z=\int(x^2+y^2)\dd m=\int x^2\dd m+\int y^2\dd m=I_y+I_x.
\end{equation}
这就是垂直轴定理。它不要求原点必须是质心，但要求三轴共点、两两垂直及物体为平面薄片，
不能直接套用于有厚度的三维物体。

\subsection{角动量守恒与陀螺进动}
系统总角动量守恒，不要求每个部分的角动量分别守恒。收拢双臂使转动惯量减小，
在外力矩可忽略且绕同一主轴的条件下，$I\omega=L$ 不变而转速增加。
此时 $E_k=L^2/(2I)$ 反而增加，能量来自肌肉做功；角动量守恒不等于机械能守恒。
图~\ref{ch09:fig:cosmic}的飓风与天体吸积结构说明旋转在不同尺度都很常见。
分析它们仍要检查外力矩、摩擦、角动量输运及物质交换，不能仅凭旋涡外形断言整个系统严格守恒。
\historyfigure[0.78\linewidth]{cosmic-angular-momentum-redraw.png}{飓风与天体吸积结构中的旋转\label{ch09:fig:cosmic}}

图~\ref{ch09:fig:gyro}中，对称陀螺固定支点为O，质心为C，$OC=r$，质量为M；
自转轴对竖直方向的倾角为$\theta$，高速自转角速度为 $\omega_0$。若自转角动量占主导，
$L\simeq I_0\omega_0$ 沿自转轴。重力对O的力矩大小 $Mgr\sin\theta$，垂直于
$\overrightarrow{OC}$ 与重力所在平面；支点反力通过O，不贡献对O的力矩。
\gyroDiagram
稳定进动时$\theta$近似不变，$\vect L$ 的端点沿水平圆周移动，圆半径为 $L\sin\theta$。
两根等长矢量 $\vect L_i,\vect L_f$ 的差连接其端点，若转过小方位角 $\dd\phi$，
\begin{equation}
 |\dd\vect L|=L\sin\theta\,\dd\phi,
 \qquad \left|\frac{\dd\vect L}{\dd t}\right|=L\sin\theta\,\omega_p.
\end{equation}
与重力矩相等，且 $\sin\theta\ne0$，得到
\begin{equation}
 Mgr\sin\theta=I_0\omega_0\sin\theta\,\omega_p,
 \qquad \boxed{\omega_p=\frac{Mgr}{I_0\omega_0}}.
\end{equation}
进动主要改变L的方向。上述式子忽略了进动本身对角动量的贡献，并要求自转足够快、倾角稳定、
支点固定；自转较慢时可能出现明显章动，即倾角随时间振荡。$\theta$=0时轴没有可辨认的绕竖直进动圆，
不能用约去零因子的步骤推断可观测的进动。式中r始终是支点到质心的距离，不是陀螺轮体半径。


\Needspace{0.34\textheight}
\begin{origexamplebox}{例题 14.1【统统算出来】}
已知 $\overrightarrow{AB}=(1,1,1)$，$\overrightarrow{AC}=(1,1,0)$，求：
\begin{enumerate}[label=(\arabic*),leftmargin=2.2em]
  \item $\overrightarrow{AB}\times\overrightarrow{AC}$；
  \item $\overrightarrow{AB}\cdot\overrightarrow{AC}$；
  \item $\overrightarrow{AB}$ 与 $\overrightarrow{AC}$ 夹角的正弦；
  \item $\overrightarrow{AB}$ 与 $\overrightarrow{AC}$ 夹角的余弦；
  \item 验证 $\sin^2\theta+\cos^2\theta=1$。
\end{enumerate}
\end{origexamplebox}

\begin{xesolution}
记 $\vect a=(1,1,1)$、$\vect b=(1,1,0)$。直接按分量计算：
\[
 \vect a\times\vect b
 =\begin{vmatrix}
   \vect i&\vect j&\vect k\\
   1&1&1\\
   1&1&0
  \end{vmatrix}
 =\boxed{(-1,1,0)}.
\]
而
\[
 \vect a\cdot\vect b=1+1+0=\boxed2,
 \qquad
 |\vect a|=\sqrt3,\quad |\vect b|=\sqrt2.
\]
由内、外积的几何意义，
\[
 \sin\theta=\frac{|\vect a\times\vect b|}{|\vect a||\vect b|}
 =\frac{\sqrt2}{\sqrt6}=\boxed{\frac1{\sqrt3}},
\]
\[
 \cos\theta=\frac{\vect a\cdot\vect b}{|\vect a||\vect b|}
 =\frac2{\sqrt6}=\boxed{\sqrt{\frac23}}.
\]
最后
\[
 \sin^2\theta+\cos^2\theta=\frac13+\frac23=1,
\]
验证完毕。由于内积为正，夹角是锐角；因此这里的正弦、余弦均取正值。
\end{xesolution}

\Needspace{0.38\textheight}
\begin{origexamplebox}{例题 14.2【面积速度】}
学习角动量之前，小明定义了“面积速度”：在参考点看来，一个质点与参考点的连线在运动
过程中会扫过一定面积。按照右手螺旋方向把面积定义为面积矢量 $\vect S$，面积速度矢量即
单位时间内扫过的面积矢量。试说明面积速度矢量与角动量的关系。
\areaVelocityDiagram
\end{origexamplebox}

\begin{xesolution}
在短时间 $\dd t$ 内，质点位移为 $\dd\vect r=\vect v\dd t$。
图~\ref{ch09:fig:area}中窄三角形的有向面积为
\[
 \dd\vect S=\frac12\vect r\times\dd\vect r
 =\frac12\vect r\times\vect v\,\dd t.
\]
所以面积速度矢量
\[
 \boxed{\frac{\dd\vect S}{\dd t}
 =\frac12\vect r\times\vect v
 =\frac{\vect L}{2m}}.
\]
面积速度与角动量方向相同，大小相差常数 $2m$。于是，有心力作用下角动量守恒等价于
面积速度恒定，这正是开普勒第二定律的几何本质。
\end{xesolution}

\Needspace{0.35\textheight}
\begin{origexamplebox}{例题 14.3【继续算出来】}
小明研究刚体的角动量，需要先计算转动惯量。以下刚体质量均匀分布，质量均为 $m$。求：
\begin{enumerate}[label=(\arabic*),leftmargin=2.2em]
  \item 长为 $l$ 的细杆绕其中垂线的转动惯量；
  \item 半径为 $R$ 的圆盘绕一条直径的转动惯量；
  \item 半径为 $R$ 的空心球壳绕一条直径的转动惯量；
  \item 半径为 $R$ 的实心球体绕一条直径的转动惯量。
\end{enumerate}
\end{origexamplebox}

\begin{xesolution}
四问都从 $I=\int r_\perp^2\dd m$ 出发；$r_\perp$ 必须是到所求转轴的垂距，而不一定是到原点的距离。

\textbf{(1) 细杆。}\quad 取杆中点为原点，线密度 $\lambda=m/l$：
\[
 I=\lambda\int_{-l/2}^{l/2}x^2\dd x
 =\frac ml\left[\frac{x^3}{3}\right]_{-l/2}^{l/2}
 =\frac ml\frac{l^3}{12}=\boxed{\frac1{12}ml^2}.
\]

\textbf{(2) 圆盘。}\quad 先算垂直盘面过圆心轴的转动惯量。面密度
$\sigma=m/(\pi R^2)$，半径为 $r$ 的圆环质量 $\dd m=2\pi r\sigma\dd r$，故
\[
 I_z=\int_0^Rr^2(2\pi r\sigma\dd r)=\frac12mR^2.
\]
对平面薄片使用垂直轴定理 $I_z=I_x+I_y$，又由对称性 $I_x=I_y$，所以绕直径
\[
 \boxed{I_x=\frac14mR^2}.
\]

\textbf{(3) 空心球壳。}\quad 这里指厚度可忽略的均匀薄球壳。取z为所求直径轴，
极角$\theta$从正z方向量起。图~\ref{ch09:fig:sphere-band}中球带的环半径为 $R\sin\theta$，
沿球面的带宽为 $R\dd\theta$，二者相乘再乘 $2\pi$ 得球带面积；它不是投影到水平面后的面积。
\sphereBandDiagram
面密度 $\sigma=m/(4\pi R^2)$，因而
\[
 \dd m=2\pi R^2\sigma\sin\theta\dd\theta,
 \qquad r_\perp=R\sin\theta.
\]
于是
\[
 I=2\pi\sigma R^4\int_0^\pi\sin^3\theta\dd\theta.
\]
令 $u=\cos\theta$，$\dd u=-\sin\theta\dd\theta$，上下限由1变为−1：
\[
 \int_0^\pi\sin^3\theta\dd\theta
 =\int_{-1}^{1}(1-u^2)\dd u
 =\left[u-\frac{u^3}{3}\right]_{-1}^{1}=\frac43.
\]
代入面密度，$I=2\pi[m/(4\pi R^2)]R^4(4/3)=\boxed{2mR^2/3}$。

\textbf{(4) 实心球。}\quad 把球体看成半径为 $r$、厚度 $\dd r$ 的球壳。设体密度
$\rho=3m/(4\pi R^3)$，球壳质量为 $\dd m=4\pi r^2\rho\dd r$，其转动惯量为
$\dd I=(2/3)r^2\dd m$，故
\[
 I=\frac{8\pi\rho}{3}\int_0^Rr^4\dd r
 =\frac{8\pi\rho R^5}{15}
 =\boxed{\frac25mR^2}.
\]
四个结果都有质量乘长度平方的量纲；同质量、同外半径下，实心球的质量更靠近轴，
其惯量小于薄球壳，$2/5<2/3$ 与几何分布一致。
\end{xesolution}

\Needspace{0.42\textheight}
\begin{origexamplebox}{例题 14.4【小明的疑惑】}
质量为 $m$ 的质点被长为 $l$ 的绳子拉着，在水平面内做圆周运动，绳子与竖直方向夹角
$\theta$。
\begin{enumerate}[label=(\arabic*),leftmargin=2.2em]
  \item 小球的速度大小为多大？
  \item 小球旋转的角速度 $\vect\omega$ 的大小和方向？
  \item 相对于悬挂点，小球角动量 $\vect L$ 的方向？
  \item 角动量与角速度似乎不在同一方向上，这样对吗？
  \item 验证这个运动是否满足角动量定理。
\end{enumerate}
\conicalPendulumDiagram
\end{origexamplebox}

\begin{xesolution}
取 $0<\theta<\pi/2$，绳不可伸长且小球可视为质点。图~\ref{ch09:fig:conical}中圆周半径
$r=l\sin\theta$。绳的张力 $T$ 的竖直分量平衡重力，水平分量提供向心力：
\[
 T\cos\theta=mg,
 \qquad
 T\sin\theta=\frac{mv^2}{l\sin\theta}.
\]
相除得
\[
 \boxed{v=\sqrt{gl\,\frac{\sin^2\theta}{\cos\theta}}}.
\]
小球绕竖直轴转动，故
\[
 \boxed{\omega=\frac{v}{l\sin\theta}
 =\sqrt{\frac{g}{l\cos\theta}}},
\]
$\vect\omega$ 沿竖直轴，方向由圆周运动方向按右手螺旋法则确定。

再看相对于悬点 $O$ 的角动量。与图一致，取 $+x$ 向右、$+z$ 向上、$+y$ 入纸面，
小球此刻在O的右下方，速度沿 $+y$，因此 $\vect\omega$ 向上。此时
\[
 \vect r=(l\sin\theta,0,-l\cos\theta),\qquad
 \vect v=(0,v,0).
\]
所以
\[
 \boxed{\vect L=m\vect r\times\vect v
 =(mlv\cos\theta,0,mlv\sin\theta)}.
\]
$\vect L$ 与绳、竖直轴共面，且垂直于绳；它一般并不与竖直的 $\vect\omega$ 平行。
这并不矛盾：只有绕对称主轴定轴转动时才能简单写成 $\vect L=I\vect\omega$，质点作圆锥摆
并不满足该条件。

最后验证角动量定理。张力沿 $\vect r$，对悬点力矩为零；只有重力产生力矩：
\[
 \vect M_O=\vect r\times m\vect g=(0,mgl\sin\theta,0).
\]
上面的 $\vect L$ 随小球以角速度 $\omega$ 绕竖直轴转动。其水平分量大小为
$L_h=mlv\cos\theta$，因而
\[
 \left|\frac{\dd\vect L}{\dd t}\right|
 =\omega L_h
 =\frac{v}{l\sin\theta}(mlv\cos\theta)
 =\frac{mv^2\cos\theta}{\sin\theta}
 =mgl\sin\theta,
\]
其方向为 $+y$：绕 $+z$ 转动使L的正x水平分量转向 $+y$。
也可直接用 $\dot{\vect L}=\vect\omega\times\vect L$
得到 $(0,\omega mlv\cos\theta,0)$，故大小和方向都与 $\vect M_O$ 一致。
当$\theta$趋于零时v趋于零，轨道缩为一点；角速度公式的有限极限不代表静止点仍有可见的圆周轨迹。
\end{xesolution}

\Needspace{0.38\textheight}
\begin{origexamplebox}{例题 14.5【哑铃落地】}
两端各为质量m质点的哑铃由长 $2l$ 的轻杆连接，以速度 $v_0$ 竖直向下纯平动，杆与地面
夹角为 $0\le\theta\le\pi/2$。下端A首先碰到沙地，并在极短碰撞中粘住、固定在该点；
另一端B不同时碰地。忽略碰撞期间重力冲量。用角动量方法求碰后瞬间角速度、B的速度及机械能损失。
这也是把碰撞与能量方法用于哑铃落地模型的另一种解法。
\dumbbellGroundDiagram
\end{origexamplebox}

\begin{xesolution}
图~\ref{ch09:fig:dumbbell-ground}标出了A、B及碰前速度。对固定接触点A取矩，
沙地冲量的力矩为零，故碰撞前后关于A的角动量守恒。取顺时针为正，B的动量产生正角动量。
$\theta$=0的情形应理解为A先接触的极限；若两端同时接触，则是另一种约束，不能套本解。

碰前上端 $B$ 的动量为 $m\vect v_0$，下端 $A$ 相对自身无力臂，所以
\[
 L_A^-=m v_0(2l\cos\theta).
\]
碰后哑铃绕 $A$ 转动，转动惯量只有上端质点贡献：
\[
 I_A=m(2l)^2=4ml^2.
\]
因此
\[
 m v_0(2l\cos\theta)=4ml^2\omega,
 \qquad
 \boxed{\omega=\frac{v_0\cos\theta}{2l}}.
\]
碰后瞬间上端速度大小为
\[
 \boxed{v_B'=2l\omega=v_0\cos\theta},
\]
方向垂直于杆。碰撞损失的机械能为
\[
 \Delta E
 =mv_0^2-\frac12I_A\omega^2
 =\boxed{mv_0^2\left(1-\frac12\cos^2\theta\right)}.
\]
这里不能对整个哑铃应用线动量守恒，因为沙地给出的外冲量不可忽略；但以碰撞点取矩后，
该冲量的力矩冲量为零，角动量守恒恰好可用。
\end{xesolution}

\Needspace{0.40\textheight}
\begin{origexamplebox}{例题 14.6【哑铃的碰撞】}
两个相同哑铃进行完全弹性碰撞：一个初始静止，另一个以速度 $v_0$ 撞来。左侧哑铃的右端
重物与右侧哑铃的左端重物发生对心正碰，且碰撞线垂直于两根轻杆。每个端点重物的质量均为
$m$。求：
\begin{enumerate}[label=(\arabic*),leftmargin=2.2em]
  \item 碰后四个重物的运动速度；
  \item 计算接触两端点的恢复系数。
\end{enumerate}
\dumbbellCollisionDiagram
\end{origexamplebox}

\begin{xesolution}
令每根杆长为 $2l$。只讨论沿碰撞线方向的速度；图~\ref{ch09:fig:dumbbell-collision}画的是碰撞前，
运动哑铃纯平动，四个端点依次记为
$A,B,C,D$，其中 $B,C$ 接触。设碰撞冲量大小为 $J$，方向使 $B$ 减速、$C$ 加速。

对左哑铃，总质量 $2m$，质心速度改变量为 $-J/(2m)$；关于质心的角冲量为 $-Jl$，而
$I_{\rm c}=2ml^2$，故角速度改变量为 $-J/(2ml)$。于是
\[
 v_A'=v_0,\qquad v_B'=v_0-\frac{J}{m}.
\]
同理对右哑铃，
\[
 v_C'=\frac{J}{m},\qquad v_D'=0.
\]
这组结果清楚地显示：冲量只改变两个接触端点的瞬时速度，远端点在碰后瞬间仍保持原速度。

完全弹性碰撞还要求系统动能守恒：
\[
 \frac12m(v_0^2+v_0^2)
 =\frac12m\left[v_0^2+\left(v_0-\frac Jm\right)^2
 +\left(\frac Jm\right)^2+0^2\right].
\]
除去无碰撞的 $J=0$ 解，得到 $J=mv_0$，因而
\[
 \boxed{v_A'=v_0,\quad v_B'=0,\quad v_C'=v_0,\quad v_D'=0}.
\]
接触点的接近速度为 $v_B-v_C=v_0$，分离速度为 $v_C'-v_B'=v_0$，所以
\[
 \boxed{e=\frac{v_C'-v_B'}{v_B-v_C}=1}.
\]
这也与题设“完全弹性碰撞”一致。
\end{xesolution}

\Needspace{0.38\textheight}
\begin{origexamplebox}{例题 14.7【跷跷板跳高】}
小明想跳得更高，以便扣篮。他把两个轻质、等臂跷跷板顺次放在篮球场上。A 同学质量
$90\,\mathrm{kg}$，B 同学质量 $80\,\mathrm{kg}$，小明质量 $60\,\mathrm{kg}$。A 从
$50\,\mathrm{cm}$ 高处落在第一块跷跷板上，把 B 弹起；B 随后落在第二块跷跷板上，把
小明弹起。所有碰撞均为弹性碰撞。按短时、等臂的理想冲击模型，只计算各人的竖直运动：
跷跷板在碰撞位置近于水平，B两次接触点等高，忽略转移到第二装置所需的水平运动对竖直运动的影响，
两次弹射后人与板均可分离。小明相对于起跳位置的重心能上升多高？
\seesawDiagram
\end{origexamplebox}

\begin{xesolution}
A 落到跷跷板时的速率为
\[
 u_A=\sqrt{2gh_0},\qquad h_0=0.50\,\mathrm m.
\]
图~\ref{ch09:fig:seesaw}用两个面板表示两次冲击，面板间的箭头只表示过程先后，不是飞行轨迹。
设第一次接触中板给A、B的竖直向上冲量大小分别为 $J_A,J_B$。对轻质板的固定支点取冲量矩，
等臂且板的转动惯量忽略，故 $J_A=J_B=J$。令A碰后向下速度为 $u_A'$、B向上速度为 $v_B$：
\[
 m_A(u_A'-u_A)=-J,\qquad m_Bv_B=J.
\]
消去J得 $m_Au_A=m_Au_A'+m_Bv_B$。完全弹性的理想传递还给出
\[
 \frac12m_Au_A^2=\frac12m_Au_A'^2+\frac12m_Bv_B^2.
\]
由第一式 $m_A(u_A-u_A')=m_Bv_B$，代入第二式并除去无传递的零解，得到
$u_A+u_A'=v_B$。再与动量形式的关系联立，第一次弹射速度为
\[
 v_B=\frac{2m_A}{m_A+m_B}u_A.
\]
B 上升后再落回同一高度，撞击第二块跷跷板前速率仍为 $v_B$。第二次弹性碰撞后，小明的
起跳速度为
\[
 v_C=\frac{2m_B}{m_B+m_C}v_B
 =\frac{2m_B}{m_B+m_C}\frac{2m_A}{m_A+m_B}u_A.
\]
故小明相对于起跳位置的上升高度
\begin{align*}
 h&=\frac{v_C^2}{2g}
 =h_0\left(
 \frac{2m_A}{m_A+m_B}
 \frac{2m_B}{m_B+m_C}
 \right)^2\\
 &=0.50\left(\frac{180}{170}\cdot\frac{160}{140}\right)^2\mathrm m
 \approx\boxed{0.732\,\mathrm m}.
\end{align*}
结果是在题设理想传递与等高条件下的重心上升量。轻板两端的瞬时几何约束不意味着A、B分离后的
自由速度相等；正是冲量关系与弹性条件共同给出结果。若B第二次落点低于第一次起跳点 $\Delta h$，
应把第二次入射速率改为 $\sqrt{v_B^2+2g\Delta h}$，不能仍用 $v_B$。
式中速度与 $\sqrt{gh_0}$ 同量纲，全部质量同比放大时高度不变；实际耗散、板质量和主动蹬跳须另建模型。
\end{xesolution}

\Needspace{0.44\textheight}
\begin{origexamplebox}{例题 14.8【漏斗与黄豆】}
一个光滑漏斗内部为圆锥面，轴截面是正三角形。小明把一粒黄豆放在漏斗内，给它合适的
水平初速度 $v_0$，黄豆在漏斗中作匀速圆周运动。
\begin{enumerate}[label=(\arabic*),leftmargin=2.2em]
  \item 黄豆转圈位置与圆心的距离为多大？
  \item 若把初速度加倍，且黄豆不会飞出或漏下，最高、最低分别能达到多高？
  \item 若把初速度减半，且黄豆不会飞出或漏下，最高、最低分别能达到多高？
\end{enumerate}
\funnelDiagram
\end{origexamplebox}

\begin{xesolution}
图~\ref{ch09:fig:funnel}以O为圆锥顶点，h竖直向上量；P处的初速度垂直轴截面，不能沿半径。
圆锥轴截面是正三角形，所以母线与竖直轴夹角 $\theta=30^\circ$，半径与高度满足
$r=h\tan\theta=h/\sqrt3$。

\textbf{(1)} 黄豆作水平匀速圆周运动。支持力 $N$ 的竖直分量平衡重力，水平分量提供向心力：
\[
 N\sin\theta=mg,
 \qquad
 N\cos\theta=\frac{mv_0^2}{r_0}.
\]
所以
\[
 \frac{v_0^2}{r_0}=g\cot\theta,
 \qquad r_0=\frac{v_0^2\tan\theta}{g}=\boxed{\frac{v_0^2}{\sqrt3g}},
 \qquad
 h_0=r_0\cot\theta=\boxed{\frac{v_0^2}{g}}.
\]

\textbf{(2)(3)} 改变初速度后，黄豆不再作定高圆周运动，但关于圆锥轴的角动量和机械能仍守恒。
令新初速度为 $u=kv_0$，其中加倍时 $k=2$，减半时 $k=1/2$。初态位于 $r_0$，在最高、
最低回归点处径向速度为零。轴向角动量守恒给出
\[
 mur_0=mvr,
 \qquad v=\frac{ur_0}{r}.
\]
机械能守恒，且 $h=\sqrt3r$：
\[
 \frac12mu^2+mg\sqrt3r_0
 =\frac12m\frac{u^2r_0^2}{r^2}+mg\sqrt3r.
\]
令 $x=r/r_0$，再用原匀速圆周条件 $v_0^2=\sqrt3gr_0=gh_0$，把能量等式除以 $mv_0^2$：
\[
 \frac{k^2}{2}+1=\frac{k^2}{2x^2}+x.
\]
乘以 $2x^2$ 并移项，化为
\[
 2x^3-(k^2+2)x^2+k^2
 =(x-1)\bigl(2x^2-k^2x-k^2\bigr)=0.
\]
$x=1$ 是改变速度的瞬间所在回归点，另一个正根为
\[
 x_* =\frac{k^2+k\sqrt{k^2+8}}{4}.
\]
因高度与半径成正比，$h_*=x_*h_0$。

当 $k=2$ 时，$x_*=1+\sqrt3>1$，初始位置是最低点：
\[
 \boxed{h_{\min}=h_0=\frac{v_0^2}{g}},
 \qquad
 \boxed{h_{\max}=(1+\sqrt3)h_0
 =\frac{(1+\sqrt3)v_0^2}{g}}.
\]
\end{xesolution}

\begin{solutioncontinuation}
当 $k=1/2$ 时，$x_*=(1+\sqrt{33})/16<1$，初始位置是最高点：
\[
 \boxed{h_{\max}=h_0=\frac{v_0^2}{g}},
 \qquad
 \boxed{h_{\min}=\frac{1+\sqrt{33}}{16}h_0
 =\frac{(1+\sqrt{33})v_0^2}{16g}}.
\]
图中的 $h$ 是从圆锥顶点所在水平面量起的竖直高度；两种结论都需要满足题设“不飞出、
不漏下”的边界条件。若k=1，方程化为 $(x-1)^2(2x+1)=0$，唯一正回归半径仍为 $r_0$，
恢复原匀速圆周；这也是对全部系数的检验。各高度都与 $v_0^2/g$ 同量纲。
光滑圆锥面的法向方程还给出 $N=m(g\sin\theta+v_\varphi^2\cos\theta/r)>0$，
说明在r为正的这段运动中无需负支持力；实际漏斗的口缘与底孔仍须按题设检查。
\end{solutioncontinuation}
